% Tapis de Quillen, with Quillen's offprint and preprint (cote 162-5, typed part): archive cote 162-5, site page 14 (PDF page 15).
% typescript. Draft compared: none.
% AI-transcribed by Claude Opus 5.5 (claude-opus-5-5, adaptive thinking, effort high), Message Batches API, batch t162_full (msgbatch_01UAh3zyWPd454AeDkhnmQ33), 2026-11-13; not yet checked by eye.
\scanpage{162-5/14}
\mspage{13}

g manifestement se prolonge en $\overline{g}$ : $\overline{T}\to X\times Y$ . Cela nous ramène à prouv
\inserted{le composé}
que \hand{(}$\unsure{\overline{T}'}\xrightarrow{\hand{\alpha}}\unsure{\overline{Z}'_i}\xrightarrow{\hand{\unsure{\alpha_i}}}\overline{T}'$\hand{)} ne dépend pas de i=0,1. Ecrivons X au lieu de $\overline{T}'$,

considérons le diagramme cartésien d'inclusions
\diagram{\begin{tikzcd}
X_s \arrow[r, "\hand{\alpha_s}"] \arrow[d, "\hand{\alpha_s}"'] & X \arrow[d, "\hand{\unsure{\Gamma_{\underline{pr}_1(x)}}}"] \\
X \arrow[r, "i_s"] & S^1\times X
\end{tikzcd}}
\[ , \]

on a alors $\hand{\alpha}^{\hand{!}}_s\,\hand{\alpha}_{s*} = \hand{\Gamma}^{\hand{!}}i_{s\hand{\unsure{*}}}$ , qui ne dépend pas de s puisque $i_s$ n'en

dépend pas à homotopie près. - Le reste est sans doute l'AQT.

Sur la strcture des B.(X). Ce sont tous des algèbres graduées

sur B.(pt)=\hand{$\Lambda$} , qui est une algèbre graduée à degrés positifs, réduit
\margin{\illegible{1}}
à $\underline{F}_2$ en degré\xout{z} 0. C'est d'ailleurs une algèbre de polynômes connue\inserted{(à une infinité de générateurs \unsure{$X_i$} \illegible{2})}
\inserted[typed]{d'anneauxpgradués}
On a un homomorphisme évident B.(X) $\to H^{\hand{\unsure{B}}}(X)$, nul en degrés \hand{$>$}0, qui
\margin{NB $H^*(X)$ $(=H^*(X,\underline{F}_2))$}
à tout Z sur X associe l'image de 1 \hand{$\in$} $H^0(Z)$ par l'homomorphisme de Gys
\margin{\unsure{en fait} $H(X)\simeq B.(X)/\Lambda^+B.(X)$}
associé à Z$\to$ X. Cet homomorphisme par Thom est surjectif (toute cla

se de cohomologie provient d'une image de variété compacte....); \xout{xxte}

On ppouve même qu'il existe un relèvement fonctoriel en X
\margin{?}
$H^\bullet$(X) $\to$ B.(X), \hand{(}compatible avec multiplication\hand{)}\inserted{il est d'ailleurs \unsure{unique s'il existe}} \xout{xxx}tel que l'on ait:

\unsure{a)} L'homomorphisme correspondant $H^\bullet(X)\otimes_{\underline{F}_2}\Lambda\simeq$ B.(X) est un iso
\margin{?}
morphisme (ce sera vrai pour tout relèvement individuellement ...), e

b) \hand{(}\deleted{Le\xout{xx}}\inserted{\unsure{quelconque}} relèvement est compatible aussi avec les hom. de Gysin

Ainsi, la catégorie $\underline{B}$ admet un foncteur pleinement fidèle dans la

catégorie des modules gradués libres de type fini sur \hand{$\Lambda$}, les modules

de type fini qu'on obtient ainsi ayant des générateurs à degrés \hand{$<$} 0

satisfaisant (pour X connexe i.e. quand il y a un seul générateur de
\margin{i.e. l'\unsure{enveloppe} \unsure{``Karoubienne''}}
degré zéro\hand{)} la dualité de Poincaré. Quand on prend la catégorie \uline{quas
\uline{abélienne}\inserted{\unsure{B}}-associée à $\underline{B}$, on trouve donc tous les modules gradués libre

de type fini à générateurs \xout{négatifsx} en degrés \hand{$\leq$} 0 seulement. C'est u

\margin{14}
